% Image credits for the photographs and AI illustrations of Book 4 (University, Year 2), Indonesian edition.
% Machine-readable license records live in images/book4/CREDITS.md; keep this
% file in sync with frontmatter/image-credits-book4.tex (one \item per item there).
% Every credit, author, dataset and licence is kept exactly as in the English file.
\chapter*{Kredit Gambar}

Gambar, peta dan grafik dalam buku ini bersifat orisinal; peta dan grafik yang
dibuat dari data yang diterbitkan menyebutkan sumbernya dalam keterangan gambar
atau dalam daftar ini. Foto digunakan di bawah lisensi bebasnya masing-masing,
dengan terima kasih kepada para penulisnya; tautan sumber lengkap dan URL lisensi
untuk setiap foto tercantum dalam \texttt{images/book4/CREDITS.md} di repositori buku.

\bigskip
Selain foto dan gambar orisinal, beberapa figur adalah ilustrasi yang dihasilkan AI,
dibuat dengan pembuatan gambar OpenAI dari prompt yang ditulis oleh editor buku, dan
ditinjau keakuratan ilmiahnya sebelum dimasukkan. Gambar-gambar ini hanya menampilkan
adegan umum, tidak pernah tempat nyata, orang nyata, atau spesimen nyata. Prompt
lengkap untuk setiap gambar yang dihasilkan tercatat dalam
\texttt{images/book4/ai/PROMPTS.md} di repositori.

\bigskip
\noindent Foto dan data, menurut bab:
\begin{itemize}\small
  \item Bab~1 (medium antarbintang dan pembentukan bintang): cakram protoplanet di sekeliling HL Tauri, ALMA (ESO/NAOJ/NRAO), CC BY 4.0; Pilar Penciptaan di Nebula Elang, NASA, ESA dan Hubble Heritage Team (STScI/AURA), domain publik.
  \item Bab~2 (pembentukan planet dan meteorit): kondrit berkarbon Allende dan pola Widmanstätten pada meteorit besi Gibeon, James St.~John, CC BY 2.0.
  \item Bab~3 (unsur-unsur di dalam Bumi): xenolit mantel lherzolit di dalam basal, Kilbourne Hole, James St.~John, CC BY 2.0.
  \item Bab~5 (magma dan penempatannya): retas di dekat Shiprock, Dexter Perkins, CC0 (Wikimedia Commons).
  \item Bab~6 (erupsi gunung api): air mancur lava di kerucut Pu`u~`\=O`\=o, Kīlauea, J.\,D. Griggs, USGS, domain publik; kolom erupsi Gunung St.~Helens pada 18 Mei 1980, Austin Post, USGS, domain publik; kolom erupsi Gunung Pinatubo pada 12 Juni 1991, Dave Harlow, USGS, domain publik.
  \item Bab~7 (pelelehan mantel): spesimen tangan dunit, Mike Davis (GeoDIL), CC0.
  \item Bab~8 (diferensiasi magma): lapisan-lapisan kromitit Kompleks Bushveld, Kevin Walsh, CC BY 2.0; gabro berlapis intrusi Skaergaard, C. Tegner, CC BY 4.0.
  \item Bab~10 (reaksi dan fasies metamorfik): sekis biru lawsonit, Ward Creek, California, dan eklogit, Almenningen, Norwegia, James St.~John, CC BY 2.0.
  \item Bab~11 (termobarometri dan lintasan P--T--t): sekis garnet, James St.~John, CC BY 2.0; butir koesit di dalam eklogit, polarisasi silang, J.\,R. Smyth, CC BY-SA 1.0.
  \item Bab~12 (pelelehan kerak: migmatit dan granit): singkapan migmatit di Nashua, New Hampshire, Kevin M. Gill, CC BY 4.0; urat leukogranit, Anne Burgess (Geograph), CC BY-SA 2.0.
  \item Bab~13 (anatomi sabuk pegunungan: Alpen): lava bantal ofiolit Chenaillet, Lionel Duchoiselle, CC BY-SA 3.0; Matterhorn dari Zermatt, Ximonic (Simo Räsänen), CC BY-SA 4.0.
  \item Bab~14 (sabuk tumbukan dan sabuk Andes dibandingkan): Dataran Tinggi Tibet antara Nepal dan Shigatse, Jan Kranendonk, CC BY-SA 4.0; Altiplano Bolivia, Luca Galuzzi, CC BY-SA 2.5.
  \item Bab~15 (masif tua dan siklus Wilson): Siccar Point, tebing tenggara, dave souza, CC BY-SA 4.0.
  \item Bab~16 (perbandingan cekungan sedimen): rig pengeboran semi-selam di Laut Utara, Erik Christensen, CC BY-SA 3.0.
  \item Bab~17 (pulau samudra): Atol Kure dari Stasiun Antariksa Internasional, NASA, domain publik.
  \item Bab~18 (Bumi Prakambrium): komatiit spinifeks di Pyke Hill, Ontario, James St.~John, CC BY 2.0; karbonat penudung dari Formasi Nuccaleena, Australia Selatan, James St.~John, CC BY 2.0.
  \item Bab~19 (Bumi Fanerozoikum): pilar-pilar kapur Monument Rocks, Kansas, James St.~John, CC BY 2.0; lapisan Bonarelli, Picenus, CC BY-SA 4.0.
  \item Bab~20 (fosil dan paleolingkungan): \emph{Waptia} dan \emph{Marrella} dari Serpih Burgess, James St.~John, CC BY 2.0; spesimen Berlin \emph{Archaeopteryx}, H.~Raab, CC BY-SA 3.0; \emph{Zoophycos} di Sagres, Amada44, CC BY 4.0.
  \item Bab~21 (termodinamika atmosfer): landasan suatu badai guntur sel super, Stephen Corfidi, NOAA, domain publik; sonde radio dari Integrated Global Radiosonde Archive, I. Durre dan kawan-kawan, NOAA NCEI, domain publik.
  \item Bab~22 (sistem cuaca): sistem tekanan rendah di atas Islandia, NASA/GSFC, MODIS Rapid Response Team, Jacques Descloitres, domain publik; mata Badai Isabel dari Stasiun Antariksa Internasional, Ed Lu, NASA, domain publik.
  \item Bab~23 (variabilitas iklim): indeks Osilasi Atlantik Utara dan Indeks Niño Samudra, Pusat Prediksi Iklim NOAA, domain publik.
  \item Bab~24 (arsip dan proksi iklim): lapisan tahunan inti es GISP2 pada kedalaman 1855~m, Anthony Gow, CRREL dari Korps Zeni Angkatan Darat AS (NOAA Paleoclimatology), domain publik; penampang speleotem travertin dari Diamond Caverns, Kentucky, James St.~John, CC BY 2.0; tumpukan bentik LR04 dari L.\,E. Lisiecki dan M.\,E. Raymo (2005), rekaman suhu EPICA Dome C dari J. Jouzel dan kawan-kawan (2007), dan komposit \ce{CO2} Antarktika dari B. Bereiter dan kawan-kawan (2015), NOAA NCEI Paleoclimatology, domain publik.
  \item Bab~25 (gaya paksa orbital dan zaman es): unsur orbit dan insolasi dari Berger dan Loutre (1991), NOAA NCEI Paleoclimatology, domain publik.
  \item Bab~26 (iklim sepanjang waktu geologi): rekaman isotop bentik Kenozoikum CENOGRID, T. Westerhold dan kawan-kawan (2020), PANGAEA, CC BY 4.0; kompilasi \ce{CO2} Fanerozoikum dari G.\,L. Foster, D.\,L. Royer dan D.\,J. Lunt (2017), \emph{Nature Communications}, CC BY 4.0.
  \item Bab~27 (perubahan iklim antropogenik): Laporan Penilaian Keenam IPCC, Kelompok Kerja~I, data tambahan Lampiran~III, CC BY 4.0.
  \item Bab~28 (bahaya dan risiko geologi): endapan aliran debris di Caraballeda, Venezuela, Lawson Smith dari Korps Zeni Angkatan Darat AS, domain publik; masjid yang tetap berdiri setelah tsunami 2004 di Aceh, USGS, domain publik.
  \item Bab~29 (endapan bijih): Tambang Bingham Canyon dari Stasiun Antariksa Internasional, NASA (awak Ekspedisi~15), domain publik; gossan Tambang Flambeau, Wisconsin, James St.~John, CC BY 2.0; data kadar--tonase dari D.\,A. Singer, V.\,I. Berger dan B.\,C. Moring (2008), USGS, domain publik.
  \item Bab~30 (geologi energi): pembangkit listrik panas bumi di dekat Geyserville, California, Departemen Pertahanan AS, domain publik; densitas karbon dioksida dari NIST Chemistry WebBook (Basis Data Referensi Standar 69).
  \item Bab~31 (metode lapangan dan laboratorium II): augit poikilitik yang melingkupi plagioklas, Strekeisen, CC BY-SA 4.0; sayatan tipis sekis garnet--mika, Jackdann88, CC BY 3.0; mikroprobe elektron di Survei Geologi Finlandia, Kallerna, CC BY-SA 4.0; hasil silika untuk bahan acuan batuan, J.\,E. Taggart dan D.\,F. Siems (2002), USGS, domain publik.
\end{itemize}
